% ==========================================================================
% GENERATED by scripts/split_paper.py. Do not edit this file.
%
% Source: imported_papers/Koopman-JEPA-TMLR-rewrite/main.tex
% Part 9: Implementation Details
%
% Edits belong upstream, in the published paper. To change what the
% chapter includes, change imports/bounded.toml and re-run
% the splitter.
% ==========================================================================
\section{Implementation Details}
\label{app:impl}
% [orig app] DONE. Verbatim except that the temporal-masking subsection is
% deleted, the Section 4.1 pointer is renamed, and the opening sentence says
% which claim the appendix supports.
This appendix supports the claim of Section~\ref{sec:protocol} that the dynamics module is the only component that varies across the four models.

The encoder architecture is selected by the tuning protocol below. The selected configuration, shared by all four methods, consists of three stride-2 convolutional blocks (channel widths 16, 32, 64) with Group Normalization and GELU activations, mapping $128{\times}128$ RGB inputs to a $16{\times}16$ spatial feature map, followed by a two-layer MLP (bottleneck dimension 256) with Layer Normalization, producing a $D{=}256$ dimensional latent vector.

\paragraph{Hyperparameter selection.} We select hyperparameters in two stages, both under a protocol fixed in advance of the final runs that scores every configuration by 1-step probe state MSE on a held-out validation split, disjoint from both the training and the test episodes. This criterion is short-horizon and penalizes representational collapse, and it does not select for the long-horizon behavior we report, so long-horizon performance was not optimized during model selection.
\begin{description}
\item[Stage 1 (shared architecture).] We first search the encoder architecture, jointly with the optimization and loss hyperparameters (Table~\ref{tab:search_space}), on the MLP-JEPA baseline alone (32 random trials per environment), and fix the winning encoder as the shared backbone for all four methods. Fixing a single architecture is what makes the comparison controlled: the dynamics module is then the only component that varies, so a performance difference can be attributed to it rather than to encoder capacity. We deliberately tune this architecture on the baseline, the most flexible model, so that the shared backbone is the one best suited to the strongest competitor and cannot be suspected of being tailored to favor the linear Koopman operators.
\item[Stage 2 (per-method optimization).] With the encoder frozen, we then tune the optimization and loss hyperparameters of each method independently (Table~\ref{tab:search_space}), again 32 trials each. The four dynamics modules induce very different optimization landscapes, for instance the spectral radius penalty and the warm-up schedule apply to only some variants, so a single shared setting would arbitrarily disadvantage some methods. Each method is therefore evaluated at the best configuration found within its own 32-trial budget.
\end{description}
Each selected configuration is finally trained for 2000 epochs and evaluated over 10 seeds. The selected per-method values are listed in Table~\ref{tab:hyperparams}.

\input{tables/table_search_space.tex}

\input{tables/table_hyperparams.tex}

\subsection{Data Collection}
Both environments use the same protocol. Episodes are collected under a uniformly random policy, and the executed action sequence is recorded and supplied to the model at training and evaluation time. Random exploration covers a wider region of the state space than expert trajectories would and prevents overfitting to a narrow behavioral distribution. Because the environments are deterministic and the actions are observed, the future trajectory is fully determined by the initial state, and long-horizon error reflects model error rather than environment stochasticity. The remaining ambiguity is in the initial condition, since the encoder receives a single frame and cannot observe velocity. This sets a floor on achievable long-horizon accuracy that applies equally to all four dynamics modules.

\begin{figure}[t]
    \centering
    \includegraphics[width=0.64\linewidth]{paper_figures/environments_overview.png}
    \caption{Sample frames from the two evaluation environments across background conditions (left to right: clean, easy, medium, hard). \textbf{Top:} Pendulum-swingup. \textbf{Bottom:} Cartpole-swingup. Both are action-conditioned and collected under a random policy.}
    \label{fig:environments}
\end{figure}

\subsection{Training Phases}
Training proceeds in two phases.
Phase~1 ($2000$ epochs) jointly trains the encoder and dynamics module using the JEPA objective (Section~\ref{sec:variants}). We retain the checkpoint with the lowest 1-step validation probe state MSE, which is the model carried into Phase~2 and evaluated.
The prediction weight $\lambda_p$ is held at zero for a per-method warm-up period (from 0 to 100 epochs, Table~\ref{tab:hyperparams}) to let the encoder establish a meaningful representation via VICReg alone, then raised to its selected value.
Phase~2 ($500$ epochs) freezes the encoder and dynamics, training only a convolutional decoder to map latent states back to pixels.
The decoder is used solely for visualization and does not influence the dynamics evaluation.

% =============================================================================
